\verseline{}{当图尔努斯从劳伦图姆城堡举起}
\verseline{}{战争的信号，号角以沙哑的鸣声震响；}
\verseline{}{当他催动烈马，推动众人拿起兵器，}
\verseline{}{人心顷刻骚动；整个拉提乌姆同时}
\verseline{5}{在惊惶的喧嚣中结盟，凶悍的青年}
\verseline{}{狂烈起来。为首的统帅墨萨波斯与乌芬斯，}
\verseline{}{还有蔑视众神的墨泽提乌斯，从四方征集}
\verseline{}{援军，使辽阔的田野失去耕作者。}
\verseline{}{维努卢斯也被派往伟大的狄俄墨得斯之城，\footnote{维努卢斯被派往狄俄墨得斯在意大利建立的城市求援。使者要转述的是图尔努斯一方对局势的陈述：特洛伊人已经立足，埃涅阿斯自称受命为王，并有诸族归附。本卷第 15—17 行所问仍是埃涅阿斯准备怎样利用局势、希望战争得到何种结局；这些意图会比对图尔努斯或拉提努斯更清楚地显现在狄俄墨得斯本人眼前。}}
\verseline{10}{向他求援，告诉他特乌克罗斯人已在拉提乌姆立足，}
\verseline{}{埃涅阿斯率船队到来，把战败的家神}
\verseline{}{带入此地，还声称命运要求他成为国王；}
\verseline{}{许多民族正归附这个达耳达诺斯人，}
\verseline{}{他的名声正在拉提乌姆广泛传开。}
\verseline{15}{至于埃涅阿斯借这一开端谋划什么，若时势相助，}
\verseline{}{又希望战争得到怎样的结局——这一切对狄俄墨得斯本人}
\verseline{}{会比对图尔努斯王或拉提努斯王显得更清楚。}
\verseline{}{拉提乌姆各处正发生这些事。拉俄墨冬的特洛伊英雄}
\verseline{}{全都看在眼里，在汹涌的忧虑中起伏不定，}
\verseline{20}{他把迅疾的心意一时分向这里，一时分向那里，}
\verseline{}{又扯往不同方向，把一切反复思量：}
\verseline{}{如同青铜盆中水面的一片颤光，}
\verseline{}{反射太阳，或映出明亮的月轮，}
\verseline{}{飞快掠过四面宽广的屋宇，又向上}
\verseline{25}{跃入空中，击中最高屋顶的嵌板。}
\verseline{}{夜已降临，疲惫的生灵遍布大地，}
\verseline{}{无论飞鸟还是牲畜，都被深沉的睡眠掌握。}
\verseline{}{这时，父亲埃涅阿斯在河岸，在寒冷天空之下，}
\verseline{}{胸中仍受悲惨战争扰乱，}
\verseline{30}{终于躺倒，让迟来的安歇流遍四肢。}
\verseline{}{此地的神祇台伯河神亲自从秀美的河流中\footnote{台伯河神兼具地方神祇与河流人格。他在梦中以白杨、青灰麻衣和芦苇显现，先确认拉提乌姆是埃涅阿斯的家园与家神安居之地，再自称“青苍的台伯”。第65行 caput 本译取“河源”；“高耸的城邑”通常联系上游伊特鲁里亚地区的诸城。}}
\verseline{}{显现，仿佛一位老人从白杨枝叶间升起；}
\verseline{}{轻薄的青灰麻衣裹住他，}
\verseline{}{荫暗的芦苇覆在发上；随后他这样开口，}
\verseline{35}{以这些话减轻埃涅阿斯的忧虑：}
\verseline{}{“诞生于神祇之族的人啊，你从仇敌手中带回}
\verseline{}{特洛伊的城，为我们保存永恒的帕加马；}
\verseline{}{劳伦图姆的土地、拉丁的田野久已等待着你。}
\verseline{}{这里是你安定的家园，也是家神安居之地；不要离去，}
\verseline{40}{也不要被战争的威胁吓倒：众神的一切愤激}
\verseline{}{与怒意都已平息。}
\verseline{}{免得你以为这些只是睡梦捏造的幻象，}
\verseline{}{你很快会发现一头巨大的母猪，躺在岸边的圣栎下，}
\verseline{}{刚刚产下三十头幼崽；她一身洁白，}
\verseline{45}{白色的幼崽环绕乳房，她就卧在地上。}
\verseline{}{那里将是城址，是你劳苦之后确切的安息地；}
\verseline{}{从这征兆起，三十年岁月轮回之后，}
\verseline{}{阿斯卡尼乌斯将建立名号昭著的阿尔巴。\footnote{Conte 2019 将第 46 行以方括号排印；该版校勘记指出若干重要传本缺此行，而第 43—46 行重现第三卷第 390—393 行。白母猪和三十头白色幼崽重现第三卷赫勒努斯的预言；三十年后阿斯卡尼乌斯建立阿尔巴，三十头幼崽与三十年相互呼应。城名 Alba 也与 alba（“白色”）相应。}}
\verseline{}{我所预言的绝不会落空。现在听好，我要用几句话}
\verseline{50}{告诉你，怎样以胜利者之身解开眼前的危局。}
\verseline{}{阿卡狄亚人来到这片海岸，他们的族系出自帕拉斯；\footnote{这里的第一个帕拉斯是阿卡狄亚族系的祖先，帕兰忒乌姆也依其名命名；他不同于厄万德之子帕拉斯，更不同于弥涅耳瓦的称号“帕拉斯”。本卷的 Pallas 可分别指这位祖先、厄万德之子帕拉斯或女神弥涅耳瓦。}}
\verseline{}{他们追随厄万德王，伴他而来，跟从他的军旗，}
\verseline{}{选定这片地方，在山间建起一座城，}
\verseline{}{依先祖帕拉斯之名，称作帕兰忒乌姆。}
\verseline{55}{他们与拉丁民族连年交战；}
\verseline{}{把他们接入你的营地，缔结同盟。}
\verseline{}{我会亲自引你沿着河岸，循河径直前行，}
\verseline{}{使船只在桨力推动下逆流而上。}
\verseline{}{起来吧，女神之子；星辰刚开始沉落时，}
\verseline{60}{依礼向朱诺献上祈祷，以恳求的誓愿}
\verseline{}{压过她的愤怒与威胁。你胜利以后，}
\verseline{}{要以敬礼酬谢我。我就是你所看见的河流，}
\verseline{}{水势充盈，擦过两岸，穿行肥沃的耕地——}
\verseline{}{青苍的台伯，最受天空喜爱的河川。}
\verseline{65}{这里是我广大的居所；我的源头从高耸的城邑之间涌出。”}
\verseline{}{说完，河神沉入深潭，}
\verseline{}{潜向水底；黑夜和睡眠也离开埃涅阿斯。}
\verseline{}{他站起身，望着天上太阳初升的}
\verseline{}{光，依礼以合拢的双掌捧起河水，}
\verseline{70}{向高空倾出这样的声音：}
\verseline{}{“宁芙们，劳伦图姆的宁芙们，诸河都由你们而生；}
\verseline{}{还有你，台伯父神，以及你神圣的水流，}
\verseline{}{请接纳埃涅阿斯，终于使他远离危难。}
\verseline{}{无论哪片水泽、哪个源头留住了怜悯我们困苦的你，}
\verseline{75}{无论你从哪片土地涌出，最为秀美，}
\verseline{}{我都会永远以敬奉、永远以祭品颂扬你——}
\verseline{}{你是赫斯珀里亚诸水长角的统治者。}
\verseline{}{只求你临在，更近地证实你的神威。”}
\verseline{}{他说着，从船队中选出两艘双层桨船，}
\verseline{80}{配好桨手，也让同伴整备兵器。}
\verseline{}{看啊，一个突如其来、令双眼惊异的奇兆：}
\verseline{}{一头洁白的母猪，带着同样洁白的幼崽，}
\verseline{}{现身林中，卧在青绿的岸边。}
\verseline{}{虔敬的埃涅阿斯把她献给你——献给你，伟大的朱诺，\footnote{白母猪和幼崽既是定居征兆，也是献给朱诺的祭牲。埃涅阿斯依河神指示，以礼仪性的恳求应对朱诺的敌意；随后台伯河整夜平息，倒流并铺展如池塘，使船只能够逆流而上。}}
\verseline{85}{奉行祭礼，把母猪与整群幼崽置于祭坛前。}
\verseline{}{台伯河神使河水整整一夜平息；那原本高涨的河流}
\verseline{}{默默倒流，停下波浪，温顺得}
\verseline{}{如同静止的池塘与宁谧的湖泽，}
\verseline{}{将水面铺平，使船桨无须搏斗。}
\verseline{90}{于是他们在顺遂的欢声中加快已经开始的航程；}
\verseline{}{涂过松脂的船身滑过浅水；水流惊异，}
\verseline{}{从未见过这景象的树林也惊异：闪亮的盾牌}
\verseline{}{和绘彩的船只，竟在河面远远漂来。}
\verseline{}{他们挥桨不息，劳顿了一夜又一日，}
\verseline{95}{越过漫长的河曲，时而被各色树木遮蔽，}
\verseline{}{在平静的水面穿行于青绿的森林。}
\verseline{}{燃烧的太阳已经攀上天空中央，}
\verseline{}{这时他们远远看见城墙、堡垒和稀疏的}
\verseline{}{屋顶——如今罗马的威势已使那里比齐天空，}
\verseline{100}{当时厄万德所拥有的一切，却如此贫弱。\footnote{叙述者知道此处后来会成为罗马，埃涅阿斯眼前却只有厄万德的贫弱聚落与稀疏屋舍；诗中不断以未来的罗马名称称呼这些尚未成形的地点。}}
\verseline{}{他们迅速转动船首，向城中驶近。}
\verseline{}{恰逢那天，阿卡狄亚国王正在城外圣林中，\footnote{厄万德正在举行赫尔库勒斯周年祭。陌生船只出现时，帕拉斯不许众人中断祭礼，自己持枪迎向来客；确认身份后，他握住埃涅阿斯的右手，把来客迎入祭宴。}}
\verseline{}{向伟大的安菲特律翁之子和众神举行周年敬礼。}
\verseline{}{他的儿子帕拉斯就在身旁，}
\verseline{105}{还有青年中的所有首领和贫寒的元老们，}
\verseline{}{一同奉上乳香，温热的鲜血在祭坛边冒烟。}
\verseline{}{他们看见高大的船只，在浓荫之间}
\verseline{}{滑向岸边，桨手悄然伏在桨上，}
\verseline{}{都被这突然的景象惊住，离开长席}
\verseline{110}{一齐站起。勇敢的帕拉斯不许他们中断}
\verseline{}{祭礼；自己一把抓起长枪，飞身迎上，}
\verseline{}{在远处土丘上喊道：“年轻人，是什么缘由迫使你们}
\verseline{}{试探陌生的水路？你们去向何方？}
\verseline{}{你们出自什么族系、哪片家园？带来和平，还是兵器？”}
\verseline{115}{父亲埃涅阿斯从高高的船尾这样回答，}
\verseline{}{手中向前伸出带来和平的橄榄枝：}
\verseline{}{“你看见的是特洛伊所生者，还有与拉丁人为敌的兵器；}
\verseline{}{拉丁人以骄横的战争驱逼我们这些流亡者。}
\verseline{}{我们前来寻找厄万德。替我们传达这些话：}
\verseline{120}{达耳达尼亚选出的首领已经来到，请求并肩作战。”}
\verseline{}{帕拉斯被如此重大的名号震住，}
\verseline{}{说道：“无论你是谁，请下船来，当面}
\verseline{}{与我的父亲交谈；请以客人身份进入我们的家神之间。”}
\verseline{}{他以手迎接，又握紧对方的右手，久久不放。}
\verseline{125}{众人向前走入圣林，把河流留在身后。}
\verseline{}{随后埃涅阿斯以友好的言辞对国王说：}
\verseline{}{“希腊人中最优秀的人啊，时势竟愿让我向你恳求，}
\verseline{}{向你伸出缠着圣带的枝条；}
\verseline{}{你虽是达那俄斯人的首领、阿卡狄亚人，我仍不曾畏惧，}
\verseline{130}{也不曾因你在血统上与阿特柔斯二子相连而畏惧。}
\verseline{}{我自己的勇气、众神神圣的启示、}
\verseline{}{两族相亲的祖先，以及你传遍大地的名声，}
\verseline{}{把我与你连在一起，命运也驱使我自愿前来。}
\verseline{}{达耳达诺斯，伊利乌姆城最初的父亲和建立者，\footnote{埃涅阿斯以共同祖源为结盟依据：达耳达诺斯由阿特拉斯之女厄勒克特拉所生；厄万德的父祖墨丘利则由阿特拉斯另一女迈亚所生。两系由此共同上溯阿特拉斯。这套谱系使希腊人与特洛伊人的旧敌关系被更早的亲缘重新连接。}}
\verseline{135}{据希腊人所说，是阿特拉斯之女厄勒克特拉所生；}
\verseline{}{他渡海来到特乌克罗斯人中。伟大的阿特拉斯}
\verseline{}{生下厄勒克特拉；他以肩膀托住天上的星界。}
\verseline{}{你们的父祖是墨丘利，洁白的迈亚}
\verseline{}{在寒冷的库勒涅山巅怀上他，又把他生下；}
\verseline{140}{而迈亚——倘若传闻也值得相信——同样由阿特拉斯}
\verseline{}{所生，正是那位举起满天星辰的阿特拉斯。}
\verseline{}{因此两族的血脉从同一源头分开。}
\verseline{}{凭借这些，我没有派使者，也没有先以机巧}
\verseline{}{安排对你的试探；是我，是我亲自把自己}
\verseline{145}{连同自己的性命交到你面前，来到你的门前恳求。}
\verseline{}{同一个道尼亚民族正以残酷的战争}
\verseline{}{追逼你我；他们相信，只要把我们逐走，}
\verseline{}{便再没有什么能够阻止他们把整个赫斯珀里亚}
\verseline{}{置于轭下，控制上方和下方濒临这土地的海。}
\verseline{150}{请接受也给出信义：我们有经得起战争的胸膛，}
\verseline{}{有勇气，也有在行动中受过考验的青年。”}
\verseline{}{埃涅阿斯说完。厄万德早已用目光端详}
\verseline{}{说话者的脸、眼睛和整个身体。}
\verseline{}{随后简短地回答：“特乌克罗斯人中最勇敢的人啊，}
\verseline{155}{我多么欣然地接纳你、认出你！你父亲的言辞、}
\verseline{}{伟大安喀塞斯的声音和面容，我都记得！}
\verseline{}{我记得拉俄墨冬之子普里阿摩斯前往萨拉弥斯，\footnote{厄万德回忆青年时代见到普里阿摩斯与安喀塞斯。普里阿摩斯前往萨拉弥斯探望姐姐赫西俄涅，又到阿卡狄亚；安喀塞斯赠给厄万德箭袋、吕喀亚箭、金织斗篷和辔头。旧日宾主记忆由礼物延续到帕拉斯，并成为当下盟约的私人根基。}}
\verseline{}{探望他姐姐赫西俄涅的王国，}
\verseline{}{随后又来到阿卡狄亚寒冷的疆界。}
\verseline{160}{那时青春初绽，才刚刚覆上我的双颊；}
\verseline{}{我惊叹特乌克罗斯人的首领，也惊叹}
\verseline{}{拉俄墨冬的儿子本人，但安喀塞斯比所有人都高。}
\verseline{}{我年轻的心燃起热望，渴望同那人交谈，}
\verseline{}{把自己的右手接入他的右手。}
\verseline{165}{我走上前，热切地带他来到斐纽斯城下。}
\verseline{}{临行时，他送我一只醒目的箭袋、吕喀亚的箭，}
\verseline{}{还有金线交织的斗篷，}
\verseline{}{以及两副黄金辔头，如今由我的帕拉斯拥有。}
\verseline{}{所以，你所寻求的右手已经与我缔结盟约；}
\verseline{170}{明日的晨光一旦重返大地，}
\verseline{}{我便让你们欢喜地带着援军离去，并以物资相助。}
\verseline{}{在此之前，既然你们以朋友身份来到这里，}
\verseline{}{就怀着敬意，同我们庆祝这不容延后的周年祭礼，}
\verseline{}{从现在起习惯坐上盟友的餐席。”}
\verseline{175}{说完这些，他命人把撤下的食物和酒杯重新摆好，}
\verseline{}{亲自让众人坐在铺草的座席上；}
\verseline{}{又特别把埃涅阿斯接上铺着长毛狮皮的榻，}
\verseline{}{请他坐入枫木制成的王座。}
\verseline{}{随后，选出的青年与祭坛的祭司争相端来}
\verseline{180}{烤熟的牛牲内脏，又在筐中堆满}
\verseline{}{精心制成的刻瑞斯赠礼，并斟上巴克科斯的酒。}
\verseline{}{埃涅阿斯与特洛伊青年一同进食，}
\verseline{}{享用整头公牛的脊肉和净祭的内脏。}
\verseline{}{饥饿消退、进食的欲望得到满足以后，}
\verseline{185}{厄万德国王说道：“这些周年祭礼、}
\verseline{}{这些依习俗举行的宴饮和这座奉献伟大神威的祭坛，}
\verseline{}{并非由空洞、不了解古老众神的迷信}
\verseline{}{强加给我们的；特洛伊客人啊，我们曾从凶险中}
\verseline{}{获救，才举行祭礼，年年更新应得的敬奉。\footnote{厄万德把周年祭礼的起源归于赫尔库勒斯杀死卡库斯、拯救当地居民；祭礼纪念这次获救，并年年更新对赫尔库勒斯的敬奉。}}
\verseline{190}{先看那边悬空的岩壁：}
\verseline{}{巨石如何远远散落，山中的居所}
\verseline{}{如何荒废，峭壁又拖下何等巨大的崩塌。}
\verseline{}{那里曾有一座洞穴，退入深广的岩腹；}
\verseline{}{半人怪物卡库斯的狰狞形体占据其中，\footnote{卡库斯是武尔坎努斯之子、口吐黑火的半人怪物。他把赫尔库勒斯的牛倒拖入洞，使蹄印方向反转；一头被藏的母牛以低鸣暴露盗窃。赫尔库勒斯三次寻门失败，最终推倒洞顶巨岩，在黑烟与火中贴身扼死卡库斯。}}
\verseline{195}{日光从来不能抵达；土地总被新近的}
\verseline{}{屠杀烘热，傲慢的门前钉着}
\verseline{}{人的头颅，苍白地悬在阴惨腐血中。}
\verseline{}{这怪物的父亲是武尔坎努斯；他从口中}
\verseline{}{吐出父亲的黑火，拖着庞大身躯行走。}
\verseline{200}{但终于有一天，时势为我们的愿望带来了}
\verseline{}{援助和神人的降临。那最伟大的复仇者}
\verseline{}{杀死三身的革律翁，因战利品而自豪，}
\verseline{}{阿尔喀得斯胜利归来，赶着巨大的}
\verseline{}{公牛经过此处；牛群占满山谷与河岸。}
\verseline{205}{可盗贼卡库斯生性凶野，什么罪行}
\verseline{}{或诡计都要冒险一试。}
\verseline{}{他从牛栏牵走四头身躯出众的公牛，}
\verseline{}{又牵走同样数目的母牛，个个姿态优美。}
\verseline{}{为使蹄印不按正常方向留下踪迹，}
\verseline{210}{他揪住尾巴，把牛倒拖进洞穴，让足迹}
\verseline{}{反转方向，将掠物藏进昏暗的岩腹；}
\verseline{}{没有一处痕迹把搜寻的人引向洞口。\footnote{Conte 2019 将第 212 行以方括号排印；该版校勘记记载，Reeve 与 Wakefield 均主张删去此行。}}
\verseline{}{与此同时，安菲特律翁之子正要把吃饱的牛群}
\verseline{}{赶出栏圈，准备离去；}
\verseline{215}{牛只临行时纷纷低鸣，整片树林充满}
\verseline{}{哀声，它们的呼叫回荡在山丘间。}
\verseline{}{洞中一头母牛回应了同伴，在广大的岩穴下}
\verseline{}{发出低鸣；这被扣留的牛背叛了卡库斯的希望。}
\verseline{}{阿尔喀得斯的痛怒顿时被狂烈的黑胆汁点燃：}
\verseline{220}{他一把抓起兵器和布满树瘤的沉重木棒，}
\verseline{}{奔向高山陡峭、直入天际的峰顶。}
\verseline{}{就在这时，我们的先人第一次看见卡库斯畏惧，}
\verseline{}{双眼慌乱；他立刻逃跑，比东风还快，}
\verseline{}{奔向洞穴，恐惧给他的双脚添上翅膀。}
\verseline{225}{他把自己关入洞中，扯断铁链，}
\verseline{}{放下那块由铁件和父亲的技艺}
\verseline{}{高悬的巨石，以这道闩障加固门柱。}
\verseline{}{看啊，提林斯的英雄已在怒火中赶到；他探遍}
\verseline{}{每个入口，脸时而转向这里，时而转向那里，}
\verseline{230}{咬牙切齿。他怒气炽热，三次}
\verseline{}{绕遍整个阿文提努斯山，三次徒然试探}
\verseline{}{石砌的门槛，三次疲惫地坐回谷中。}
\verseline{}{有一块尖利的巨岩，从四面削断的岩层间}
\verseline{}{高高立在洞穴背上，望去极其高峻，}
\verseline{235}{正适合凶恶的鸟类筑巢。}
\verseline{}{它从山脊向左倾斜，俯压河面；英雄从右边}
\verseline{}{朝相反方向使力，摇动巨石，}
\verseline{}{把它从最深的根基拔松；随后猛然}
\verseline{}{推落。巨响震彻最高的天空，}
\verseline{240}{两岸迸开，惊骇的河水倒退。}
\verseline{}{卡库斯巨大的洞窟与王宫就这样裸露出来，}
\verseline{}{深处荫暗的岩腔完全敞开；}
\verseline{}{仿佛大地从腹中被某种力量劈裂，}
\verseline{}{揭开地下的居所，让苍白的国度显现——}
\verseline{245}{那是众神厌恶之地；从上方望见巨大的深渊，}
\verseline{}{亡魂在闯入的光线下惊惶。}
\verseline{}{卡库斯被意外的日光骤然捕获，}
\verseline{}{困在岩穴中，发出从未有过的咆哮；}
\verseline{}{阿尔喀得斯从上方以飞射的兵器压住他，召来}
\verseline{250}{一切可以使用的武器，以树枝和巨大磨石进逼。}
\verseline{}{怪物已经没有任何道路能逃离危险，}
\verseline{}{便从咽喉吐出浓烟——说来令人惊异——}
\verseline{}{以盲目的黑雾裹住整座洞窟，}
\verseline{}{夺去双眼的视线；他在岩穴深处聚拢}
\verseline{255}{冒烟的黑夜，把烈火与幽暗混在一起。}
\verseline{}{阿尔喀得斯的心再也不能忍受；他纵身}
\verseline{}{跃入火中，正扑向浓烟翻卷最甚、}
\verseline{}{巨大洞穴在乌黑雾气中沸腾之处。}
\verseline{}{他在黑暗中抓住仍徒然喷火的卡库斯，}
\verseline{260}{以双臂把他绞成一个死结，贴身扼紧，}
\verseline{}{直至双眼迸出，咽喉枯干无血。}
\verseline{}{黑暗的居所顷刻敞开，门石被掀落；}
\verseline{}{被偷走、又被盗贼矢口否认的牛只}
\verseline{}{重见天日，那具不成形的尸体也被抓住双脚拖出。}
\verseline{265}{人们久久观看，心意仍不能满足：怪物可怖的双眼、}
\verseline{}{面孔、鬃毛般硬刺覆盖的胸膛，}
\verseline{}{还有咽喉中已经熄灭的火焰，全都摆在眼前。}
\verseline{}{从此，这项敬礼年年举行，后代欣然守住此日；\footnote{波提提乌斯被写成祭礼的创立者，皮纳里乌斯家族则守护赫尔库勒斯圣仪；圣林中的祭坛后来以“至大祭坛”著称。白杨叶一面深、一面浅，因而称“双色”；它与酒杯、奠酒和共同呼神组成祭礼的现时动作。}}
\verseline{}{波提提乌斯最先创立祭礼，皮纳里乌斯家族}
\verseline{270}{则成为赫尔库勒斯圣仪的守护者。}
\verseline{}{波提提乌斯在圣林中立下这座祭坛；}
\verseline{}{我们将永远称它为至大祭坛，它也将永远至大。}
\verseline{}{所以来吧，年轻人，为敬奉如此显赫的功业，}
\verseline{}{以枝叶环绕头发，右手举起酒杯，}
\verseline{275}{呼求我们共同的神，欣然献上美酒。”}
\verseline{}{他说完，双色白杨——赫尔库勒斯之树——以荫影}
\verseline{}{遮住众人的头，枝叶交缠着垂落；}
\verseline{}{神圣的酒杯盛满美酒，交到众人右手。他们迅速而欢喜地}
\verseline{}{在长席上奠酒，向众神祈祷。}
\verseline{280}{与此同时，天空西斜，暮星已经更贴近奥林匹斯；}
\verseline{}{祭司们走来，波提提乌斯行在最前，}
\verseline{}{依习俗以兽皮束身，手中举着火焰。}
\verseline{}{他们重新摆开宴席，端来第二席可喜的}
\verseline{}{赠礼，以盛满食物的大盘堆满祭坛。}
\verseline{285}{随后萨利祭司来到燃烧的祭坛周围歌唱，\footnote{萨利祭司分为青年与老人两队歌唱。赞歌先列赫尔库勒斯既有功业，再把卡库斯洞穴和喷火怪物纳入其中，使发生在此地的新胜利进入公共祭歌。安菲特律翁之子、阿尔喀得斯和提林斯的英雄都是赫尔库勒斯的称谓。}}
\verseline{}{额头缠着白杨枝；}
\verseline{}{一边是青年合唱，一边是老人。他们以歌声传述}
\verseline{}{赫尔库勒斯的荣耀与功业：他最初如何亲手扼杀}
\verseline{}{继母派来的两个怪物——那一双缠身的蛇；}
\verseline{290}{又如何在战争中摧毁显赫的城邦，}
\verseline{}{特洛伊和俄卡利亚；如何在欧律斯透斯王手下，}
\verseline{}{依不公的朱诺之命，承受千般艰苦的劳役：}
\verseline{}{“不可征服者啊，你亲手杀死云中诞生、半人半马的}
\verseline{}{许莱乌斯和福卢斯，也杀死克里特的}
\verseline{295}{怪物，以及涅墨亚岩下那头巨狮。}
\verseline{}{斯提克斯的湖水在你面前战栗，奥尔库斯的守门者也战栗；}
\verseline{}{它在血腥洞穴中躺在啃剩的骨头上。}
\verseline{}{没有一种形貌能使你畏惧，就连身形耸立、}
\verseline{}{手执武器的提丰也不能；}
\verseline{300}{多头环拥的勒耳那水蛇也没有使你无计可施。}
\verseline{}{向你致敬，朱庇特真正的儿子，为众神增添的荣耀；}
\verseline{}{请你亲临，以吉祥的脚步护佑我们和你的祭礼。”}
\verseline{}{他们以歌声颂扬这些事，尤其把卡库斯的}
\verseline{}{洞穴和那喷火的怪物本人也加入歌中。}
\verseline{305}{整片树林与歌声共鸣，山丘回荡。}
\verseline{}{神圣仪式完成后，众人从那里返回城中。}
\verseline{}{国王年老体衰，缓步前行，}
\verseline{}{让同行的埃涅阿斯和儿子伴在身旁，}
\verseline{}{以不断变换的话题减轻路途的劳顿。}
\verseline{310}{埃涅阿斯惊叹，让目光自在地游遍四周，}
\verseline{}{被这些地方吸引，满怀欣喜地逐一}
\verseline{}{询问、聆听前人留下的遗迹与故事。}
\verseline{}{于是，未来罗马城堡的建立者厄万德王说道：\footnote{厄万德所讲的是拉提乌姆的地方神话史：萨图尔努斯流亡至此，聚合山居者、制定法律，以自己曾“隐藏”于此解释 Latium 之名；黄金时代之后才出现战争与占有欲。河名又被追溯到巨人台伯，取代旧名阿尔布拉。这些神话性词源把地方景观、神祇和古老统治联在一起。}}
\verseline{}{“这些树林曾由本土的法乌努斯与宁芙占据，}
\verseline{315}{还有一族从树干和坚硬橡木中诞生的人；}
\verseline{}{他们没有习俗，也没有耕作技艺，不懂给公牛套轭，}
\verseline{}{不懂积聚财物，也不懂节省所得，}
\verseline{}{只靠树木的果实与艰苦的狩猎养活自己。}
\verseline{}{萨图尔努斯最先从天上的奥林匹斯来到这里，}
\verseline{320}{躲避朱庇特的兵器，在被夺去王国后流亡。}
\verseline{}{他聚合那散居高山、尚未受教的族类，}
\verseline{}{给他们制定法律；又愿意把这片土地称作拉提乌姆，}
\verseline{}{因为他曾安然隐藏在这里的海岸。}
\verseline{}{传说中的黄金时代就在那位国王治下；}
\verseline{325}{他以安宁的和平统治众民，}
\verseline{}{直到更坏、更失去光彩的时代逐渐接替，}
\verseline{}{战争的狂烈和占有的欲望随之而来。}
\verseline{}{后来奥索尼亚人的队伍和西坎尼亚民族来到，}
\verseline{}{萨图尔努斯的土地也屡次更换名称；}
\verseline{330}{又有许多国王，还有身形庞大、性情凶厉的台伯，}
\verseline{}{意大利人后来依他的名字，把这条河称为台伯河；}
\verseline{}{古老的阿尔布拉河从此失去了真名。}
\verseline{}{我被逐出故乡，追随大海最遥远的边界，}
\verseline{}{全能的福尔图娜与不可逃脱的命运}
\verseline{335}{把我安置在这里；我母亲、宁芙卡耳门塔的可畏警示}
\verseline{}{与引导者阿波罗一道催我前来。”}
\verseline{}{话音刚落，他继续向前，指出一座祭坛，\footnote{卡耳门塔门、庇护所、卢佩耳卡尔、阿耳吉勒图姆、塔耳佩亚居地和卡庇托利乌姆都是后世罗马地名投向当下荒丘的重影。卡耳门塔被写成最早预言伟大埃涅阿斯后裔与帕兰忒乌姆声名的先知；阿耳吉勒图姆则联系客人阿耳戈斯之死的地方传说。}}
\verseline{}{又指出罗马人以卡耳门塔之名称呼的城门，}
\verseline{}{那是古来对宁芙卡耳门塔的纪念；}
\verseline{340}{她是宣示命运的预言者，最先歌唱未来伟大的}
\verseline{}{埃涅阿斯后裔和声名显赫的帕兰忒乌姆。}
\verseline{}{再往前，他指给来客看一片巨大的圣林，勇烈的罗慕路斯}
\verseline{}{后来把它设为庇护所；又在寒冷岩壁下指出卢佩耳卡尔，}
\verseline{}{依帕耳拉西亚习俗，得名于吕凯奥斯的潘。}
\verseline{345}{他还指出神圣的阿耳吉勒图姆林地，}
\verseline{}{为那地方作证，讲说客人阿耳戈斯在那里死去。}
\verseline{}{然后他领众人走向塔耳佩亚居地和卡庇托利乌姆——}
\verseline{}{如今一片金色，当时却是灌木森然的荒丘。}
\verseline{}{那时，地方严厉的神圣威慑已使惶恐的乡民}
\verseline{350}{惊惧；那时，他们已会面对林木与岩石战栗。}
\verseline{}{“这片林苑，”他说，“还有这座枝叶覆盖山顶的丘陵，}
\verseline{}{住着一位神——是哪位神，没有人确定。阿卡狄亚人}
\verseline{}{相信自己看见过朱庇特：他常以右手摇动}
\verseline{}{乌黑的神盾，召集暴雨的云。”\footnote{乡民已敬畏卡庇托利乌姆的林与岩，却不知道居住其间的是哪位神；阿卡狄亚人只相信自己见过朱庇特摇动神盾、召来风暴。随后提到的雅努斯、萨图尔努斯旧城废墟，也都位于后来罗马所在的地区。}}
\verseline{355}{“此外，你还看见两座城，墙垣已经倾倒，}
\verseline{}{是古人的遗存与记忆。}
\verseline{}{一座城堡由雅努斯父神建立，另一座由萨图尔努斯建立；}
\verseline{}{前者名为雅尼库卢姆，后者名为萨图尔尼亚。”}
\verseline{}{他们这样一路交谈，走近贫寒厄万德的\footnote{牛群在罗马广场与富丽卡里奈低鸣，贫寒王宫却以宾主尊严接纳埃涅阿斯。厄万德要求客人蔑视财富、不要轻视贫乏；给埃涅阿斯铺的卧榻只有树叶和一张利比亚母熊皮。}}
\verseline{360}{屋舍，看见牛群四处游荡，}
\verseline{}{在未来的罗马广场和富丽的卡里奈低鸣。}
\verseline{}{来到居处时，厄万德说：“胜利的阿尔喀得斯}
\verseline{}{曾跨过这道门槛，这座王宫曾接纳他。}
\verseline{}{客人啊，要敢于蔑视财富，也使自己配得上神；}
\verseline{365}{来面对这些贫乏的事物，不要心怀轻蔑。”}
\verseline{}{说完，他把身躯高大的埃涅阿斯领到}
\verseline{}{低矮屋顶之下，让他躺在铺好的寝处，}
\verseline{}{以树叶和利比亚母熊的毛皮托住身体。}
\verseline{}{黑夜坠落，以昏暗的双翼拥住大地。}
\verseline{370}{母亲维纳斯心中并非无故惊惧；}
\verseline{}{劳伦图姆的威胁和艰难的战乱使她不安，}
\verseline{}{她便在丈夫金色的婚房中向武尔坎努斯开口，}
\verseline{}{以这些话开始，把神圣的爱欲吹入言辞：}
\verseline{}{“当阿尔戈斯诸王以战争摧毁帕加马——}
\verseline{375}{那座欠下覆亡之债、即将倒在仇敌火中的城堡——}
\verseline{}{我不曾为受难者请求任何援助，也不曾请求}
\verseline{}{由你的技艺和力量铸成武器；最亲爱的丈夫，}
\verseline{}{我不愿让你和你的劳作徒然费力，}
\verseline{}{尽管我欠普里阿摩斯的儿子们许多恩情，}
\verseline{380}{也曾多次为埃涅阿斯艰苦的劳顿哭泣。}
\verseline{}{如今他依朱庇特的命令，在鲁图利人的海岸立足；}
\verseline{}{所以我如今却来恳求，向我视为神圣的}
\verseline{}{神明求取武器——母亲替儿子而求。涅柔斯之女，}
\verseline{}{提托诺斯的妻子，都曾以泪水使你屈从。}
\verseline{385}{看吧，多少民族正在结集，多少城邑关闭}
\verseline{}{城门，磨利刀剑，要毁灭我和我的族人。”}
\verseline{}{她说完，女神以雪白的双臂从两边}
\verseline{}{抱住仍在迟疑的丈夫，在柔软的拥抱中温暖他。}
\verseline{}{他突然接纳了熟悉的火焰；惯常的热流}
\verseline{390}{进入骨髓，摇撼每根骨骼而奔走，}
\verseline{}{如同雷鸣时，一道灼亮的火隙突然裂开，}
\verseline{}{以闪光贯穿云层。}
\verseline{}{妻子发觉了，因计策奏效而欢喜，也深知自己的美貌。}
\verseline{}{于是，父神仍受永恒的爱所缚，开口说道：}
\verseline{395}{“为何从如此深远处寻找理由？女神啊，}
\verseline{}{你对我的信赖去了哪里？若当年你也有同样的忧心，}
\verseline{}{那时我本来也可以为特乌克罗斯人铸造武器；}
\verseline{}{全能的父神和命运都未曾禁止特洛伊}
\verseline{}{再屹立十年，也未禁止普里阿摩斯再活十年。\footnote{维纳斯以母亲身份求取武器，也明确运用拥抱、美貌和计策。武尔坎努斯指出：若她当年同样忧虑，他本可武装特洛伊人，特洛伊和普里阿摩斯本可再存十年。}}
\verseline{400}{而现在，若你准备战争，心意既已如此，}
\verseline{}{凡我的技艺能够尽心许诺的，}
\verseline{}{凡以铁、以熔化的金银合金能够制成之物，}
\verseline{}{凡烈火与风息能够完成之事——别再恳求了，}
\verseline{}{仿佛不相信你自己的力量。”说完这些话，}
\verseline{405}{他给了她所盼望的拥抱，偎入妻子的怀中，}
\verseline{}{任安宁的睡眠流遍四肢。}
\verseline{}{此后，夜程已经过半，最初一阵安睡}
\verseline{}{已把睡意驱尽；正如一位妇人起身——}
\verseline{}{命运要她靠纺锤和弥涅耳瓦微薄的手艺\footnote{本卷第 407—415 行以贫妇夜间纺织比拟武尔坎努斯起身工作。“纺锤”和“微薄的弥涅耳瓦手艺”指靠纺织维生的劳动；她重燃灰烬，令婢女继续长工，以维持婚床与抚养孩子。这个明喻先写凡人妇女夜间劳作，再转回武尔坎努斯起身前往锻炉。}}
\verseline{410}{维持生活——重新挑起灰烬中沉睡的火；}
\verseline{}{她把黑夜加进劳作，让婢女们在灯光下}
\verseline{}{继续漫长的纺织，以便守住丈夫}
\verseline{}{贞洁的婚床，也把幼小的儿女抚养长大；}
\verseline{}{火之主同样不肯在这时懈怠，}
\verseline{415}{从柔软的寝具上起身，走向工匠的劳作。}
\verseline{}{在西坎尼亚海岸与埃俄利亚的利帕拉旁，}
\verseline{}{一座岛屿耸立，冒烟的岩石高插天空；\footnote{该岛通常认作希厄拉，即今武尔卡诺岛，诗中称为 Volcania；岛下是武尔坎努斯与独眼巨人的工场。“埃特纳岩窟”联系独眼巨人的锻炉传统。布戎忒斯、斯忒洛佩斯和皮剌克蒙正铸造雷霆、玛尔斯战车与女神胸甲；雷霆由雨、云、火、南风、声响、恐惧和怒意组成。electrum 指古代所称的金银合金；埃涅阿斯之盾则由七层圆环叠合而成。}}
\verseline{}{岛下，洞穴和被独眼巨人炉火蚀空的}
\verseline{}{埃特纳岩窟轰鸣；铁砧上的重击}
\verseline{420}{在岩腔里回响，化作呻吟；卡吕布斯的钢块}
\verseline{}{在洞内嘶叫，熔炉中的火焰喘息。}
\verseline{}{那里是武尔坎努斯的居所，整片土地因他得名。}
\verseline{}{御火者就在这时从高天降到此地。}
\verseline{}{独眼巨人正在巨大的洞窟中锻打钢铁：}
\verseline{425}{布戎忒斯、斯忒洛佩斯和赤裸着肢体的皮剌克蒙。}
\verseline{}{他们手中已有一道雷霆初具形状，部分磨亮——}
\verseline{}{众神之父常把无数这样的雷霆从整个天空}
\verseline{}{投向大地；仍有一部分尚未完成。}
\verseline{}{他们已加上三道扭曲的雨、三道含水的云，}
\verseline{430}{三道赤红的火与三道长翼南风；}
\verseline{}{此刻正把可怖的闪光、轰鸣和恐惧}
\verseline{}{混入作品，也把怒火混入紧随的火焰。}
\verseline{}{另一边，他们赶制玛尔斯的战车和飞转的车轮，}
\verseline{}{战神凭它们激起男人，也激起城邦；}
\verseline{435}{还有可怖的胸甲——发怒的帕拉斯女神的武装——\footnote{此处 Palladis arma 属于弥涅耳瓦的称号“帕拉斯”；蛇鳞、纠缠的群蛇与断颈后仍转动双眼的戈耳工，都是女神胸甲上的可怖装饰。第八卷的青年帕拉斯是同名人物。}}
\verseline{}{他们争相磨亮它的蛇鳞与黄金装饰，}
\verseline{}{还有纠缠相连的群蛇；女神胸前那戈耳工，}
\verseline{}{颈项已被斩断，双眼仍在转动。}
\verseline{}{“把一切收起来，”他说，“放下已经开始的劳作，}
\verseline{440}{埃特纳的独眼巨人，把心思转向这里：}
\verseline{}{现在要为一位刚勇的男子铸造兵器。现在需要力量，}
\verseline{}{现在需要迅疾的双手，现在需要大师的一切技艺。}
\verseline{}{赶快，别再耽搁。”他没有多说；众人却立刻}
\verseline{}{一齐俯身劳作，平均分担各自的工序。}
\verseline{445}{青铜如河流淌，熔化的黄金也在流动，}
\verseline{}{伤人的钢铁在巨炉中熔化。}
\verseline{}{他们铸出一面巨盾，足以独自抵御}
\verseline{}{拉丁人的一切兵器，又以层层圆环套合七层盾面。}
\verseline{}{一些人让鼓胀的风箱吸入空气，}
\verseline{450}{再把它吐出；另一些人把嘶鸣的青铜浸入水池。}
\verseline{}{整个洞穴在架起的铁砧下呻吟；}
\verseline{}{他们的手臂以巨大力量相互应和，依节奏举起，}
\verseline{}{又用坚咬的铁钳翻转金属块。}
\verseline{}{当来自勒姆诺斯的父神在埃俄利亚海岸加紧劳作时，}
\verseline{455}{温养万物的晨光已从低矮的屋顶下唤醒厄万德，}
\verseline{}{清晨的鸟鸣也在屋檐下响起。}
\verseline{}{老人起身，把短衣穿在身上，}
\verseline{}{又以伊特鲁里亚式鞋带束住脚掌；}
\verseline{}{随后把忒革亚的剑系在身侧与肩头，}
\verseline{460}{从左边反披下垂的豹皮。}
\verseline{}{两只守门的狗也从高高的门槛前走下，}
\verseline{}{行在主人前面，伴随他的脚步。}
\verseline{}{这位英雄前去寻找客人埃涅阿斯的住处，}
\verseline{}{记得昨日的谈话和许诺过的援助。}
\verseline{465}{埃涅阿斯也同样在清晨起身。}
\verseline{}{儿子帕拉斯陪着厄万德，同伴阿卡忒斯则陪着埃涅阿斯。}
\verseline{}{两方相见，握住彼此的右手，坐在屋舍中央，}
\verseline{}{终于自在地展开不受拘束的交谈。}
\verseline{}{国王先开口：}
\verseline{470}{“特乌克罗斯人最伟大的统帅啊，只要你安然无恙，}
\verseline{}{我绝不会承认特洛伊的事业和王国已经战败。}
\verseline{}{与这般盛名相比，我们援战的力量}
\verseline{}{实在微小；一边有台伯河阻住我们，}
\verseline{}{一边有鲁图利人逼近，以兵器在城墙四周喧响。}
\verseline{475}{但我准备让辽阔的民族与汇聚众王的富足军营}
\verseline{}{同你联合——意外的机缘正显示这条生路；}
\verseline{}{命运召你来到此地，你也亲自带着自己前来。}
\verseline{}{离这里不远，有一座以古老岩石为基的城，\footnote{阿居拉是伊特鲁里亚城市凯雷的古名。厄万德把墨泽提乌斯描述为以傲慢统治和残酷武力控制此城的暴君，并详细控诉他把活人与尸体手手、口口相连的酷刑；这些都是厄万德劝埃涅阿斯结盟时对墨泽提乌斯的指控。公民起义后，墨泽提乌斯逃往鲁图利人领地，受图尔努斯庇护。}}
\verseline{}{那是阿居拉城的居地；吕底亚民族昔日}
\verseline{480}{以战争闻名，在伊特鲁里亚山脊上定居。}
\verseline{}{城市兴盛多年，后来却被墨泽提乌斯王}
\verseline{}{以傲慢的统治和残酷的兵器控制。}
\verseline{}{我何必列举不可言说的屠杀，何必列举暴君}
\verseline{}{凶野的行径？愿众神把这些留给他自己和他的族类！}
\verseline{485}{他甚至把死人的身体与活人连在一起，}
\verseline{}{手贴着手，嘴贴着嘴——}
\verseline{}{这是一种酷刑；脓血和腐液从尸身流下，}
\verseline{}{他让活人在悲惨的拥抱中受尽漫长的死亡。}
\verseline{}{疲惫的公民终于拿起武器，把仍在肆虐的暴君}
\verseline{490}{和他的宅邸包围，}
\verseline{}{斩杀他的同党，把火焰投向屋顶。}
\verseline{}{墨泽提乌斯却从屠杀中逃出，来到鲁图利人的田野，}
\verseline{}{投奔图尔努斯这位庇护者，由他的武力护卫。}
\verseline{}{于是整个伊特鲁里亚怀着正当的狂怒起身；}
\verseline{495}{他们要求立刻交战，把国王交来受刑。}
\verseline{}{埃涅阿斯，我要让你统率这成千人马。}
\verseline{}{密集的战船已经铺满整个海岸，喧声如雷，}
\verseline{}{众人要求举起军旗；一位年迈的占卜祭司却拦住他们，}
\verseline{}{吟唱命运说：‘迈俄尼亚选出的青年啊，\footnote{“迈俄尼亚青年”以吕底亚旧称指伊特鲁里亚人，呼应他们源出吕底亚的传统。占卜祭司宣布意大利人不得统领这支民族，必须选择外来领袖；塔耳孔因此把冠冕、权杖与王权标志送给厄万德。厄万德年老，帕拉斯又因萨宾母系带有本土血缘，外来统帅的位置遂转给埃涅阿斯。}}
\verseline{500}{你们是古老先民的精华，是男子的勇力；正当的悲痛}
\verseline{}{把你们推向仇敌，墨泽提乌斯又使你们燃起理所应有的怒火。}
\verseline{}{神律不容任何意大利人使这样伟大的民族屈服；}
\verseline{}{去选择外来的统帅。’伊特鲁里亚军队便停驻在}
\verseline{}{这片原野，因众神的警示而畏惧。}
\verseline{505}{塔耳孔亲自派使者到我这里，把王国的冠冕}
\verseline{}{连同权杖送来，又托付王权的标志，}
\verseline{}{要我进入军营，承受伊特鲁里亚人的王权。}
\verseline{}{但我年老体衰，寒气凝滞、岁月耗竭，}
\verseline{}{无力掌权，迟暮的力量也担不起勇烈的事业。}
\verseline{510}{我本可勉励儿子，只因他的母亲出自萨宾人，}
\verseline{}{他也从这片土地承得一部分故乡血脉。至于你，}
\verseline{}{命运宽待你的年岁和族系，众神的旨意正在召你；}
\verseline{}{前去吧，特乌克罗斯人与意大利人中最勇敢的统帅。}
\verseline{}{此外，我要把帕拉斯——我的希望和慰藉——}
\verseline{515}{交给你；让他在你这位导师身边习惯承受}
\verseline{}{军旅和玛尔斯沉重的事业，观看你的行动，}
\verseline{}{从年少之初便仰望你。}
\verseline{}{我会给他两百名阿卡狄亚骑兵，都是青年中的精锐；}
\verseline{}{帕拉斯也将以自己的名义给你同样的人数。”}
\verseline{520}{他刚刚说完，安喀塞斯之子埃涅阿斯和忠实的}
\verseline{}{阿卡忒斯都低垂着脸，久久不语，}
\verseline{}{悲沉的心里盘旋着许多艰难——若非}
\verseline{}{库忒瑞亚已在开阔的天空给出征兆。}
\verseline{}{一道闪电突然从高空抖落，伴着巨响来到，}
\verseline{525}{一切仿佛顷刻崩坠；伊特鲁里亚号角般的吼声}
\verseline{}{在天空中长鸣。}
\verseline{}{众人抬头望去；巨大的轰声一遍又一遍轰击。}
\verseline{}{他们看见兵器在云间、在一片晴朗天域中}
\verseline{}{透过澄空闪着红光，彼此撞击，轰然作响。\footnote{晴空中闪红、相互撞击的兵器尚未降落，是维纳斯预先约定的征兆。埃涅阿斯辨认出母亲的许诺，却随即预见劳伦图姆人的大屠杀、图尔努斯的惩罚，以及台伯河卷走盾牌、头盔与尸体；胜利信号与真实死亡在同一段中出现。}}
\verseline{530}{其余人心中惊愕，特洛伊英雄却}
\verseline{}{认出那声音和女神母亲的许诺。}
\verseline{}{于是他说：“客人啊，别再问了，千万不要追问}
\verseline{}{这些奇象预示怎样的遭际：奥林匹斯正在召我。}
\verseline{}{母亲女神曾预告：若战争迫近，她会发出这信号，}
\verseline{535}{还会把武尔坎努斯铸造的兵器穿过天空}
\verseline{}{送来助战。}
\verseline{}{唉，可怜的劳伦图姆人面前将有多大的屠杀！}
\verseline{}{图尔努斯，你要向我偿还何等惩罚！台伯父神，}
\verseline{}{你又要把多少盾牌、头盔和强健的尸体}
\verseline{540}{卷入水下！让他们要求列阵、撕毁盟约吧。”}
\verseline{}{说完，他从高高的王座上起身，}
\verseline{}{先在赫尔库勒斯祭坛上重新唤醒沉睡的火，}
\verseline{}{又朝昨日的炉火和微小的家神}
\verseline{}{欣然走近；众人依习俗宰杀选出的两岁绵羊，}
\verseline{545}{厄万德参与其中，特洛伊青年亦然。}
\verseline{}{此后，埃涅阿斯走向船只，重新巡视同伴，}
\verseline{}{从他们中选出勇力出众、将要跟随自己}
\verseline{}{走向战争的人；其余人则顺着下行的水势，}
\verseline{}{安缓地沿顺流漂去，}
\verseline{550}{要把事情的进展和父亲的消息带给阿斯卡尼乌斯。}
\verseline{}{特乌克罗斯人准备前往伊特鲁里亚田野，众人给他们马匹；}
\verseline{}{另牵来一匹不经抽签、专供埃涅阿斯的战马，}
\verseline{}{黄褐的狮皮从全身垂下，金色利爪灿然发亮。}
\verseline{}{消息骤然传开，飞遍小小的城邦：\footnote{这里的Fama指迅速传开的消息。城中母亲因子弟出征而加倍祈祷；塞尔维乌斯（8.557）把“玛尔斯的形象”解释为她们所想见的战争景象。这一段写母亲们的恐惧，随后转入厄万德送别帕拉斯的祈祷（558行以下）。}}
\verseline{555}{骑兵正迅速前往伊特鲁里亚国王的门庭。}
\verseline{}{母亲们在恐惧中加倍祈祷；危险越近，}
\verseline{}{恐惧越向前逼来，玛尔斯的形象已显得更加巨大。}
\verseline{}{这时父亲厄万德拥住将要离去的儿子右手，\footnote{厄鲁卢斯由母亲费罗尼亚赋予三条生命，须持三副武装、三次被击倒，最终仍被年轻的厄万德杀死。厄万德接着祈求：若帕拉斯能平安归来，他愿忍受一切；若灾祸已定，他宁愿在希望未灭、仍拥着儿子时死去。随后叙述者称这是父子的“最后一次离别”。}}
\verseline{}{久久不放，怎么也抱不够，哭着这样说道：}
\verseline{560}{“啊，但愿朱庇特把逝去的岁月还给我——}
\verseline{}{让我恢复当年模样：我曾在普莱内斯特城下击溃}
\verseline{}{最前列的军阵，胜利后焚烧成堆的盾牌，}
\verseline{}{又以这只右手把厄鲁卢斯王送入塔耳塔洛斯；}
\verseline{}{他出生时，母亲费罗尼亚给了他三条生命——}
\verseline{565}{说来可怖——因此要操持三副武装，}
\verseline{}{要三次被死亡击倒；可当时这只右手}
\verseline{}{仍夺走了他的每一条生命，三次剥下他的武装。}
\verseline{}{若我仍是那样，如今绝不会有任何力量把我从你}
\verseline{}{甜蜜的拥抱中撕开，儿子；邻近此地的墨泽提乌斯}
\verseline{570}{也绝不会骄横地欺侮我，在刀剑下造出}
\verseline{}{这样多残酷的死亡，使城中失去如此多公民。}
\verseline{}{但你们啊，天上的众神，还有你，诸神最伟大的统治者}
\verseline{}{朱庇特，我恳求你们怜悯阿卡狄亚国王，}
\verseline{}{听见一个父亲的祷告。如果你们的神意、}
\verseline{575}{如果命运要让帕拉斯平安回到我身边，}
\verseline{}{如果我能活着再见他，重新与他团聚，}
\verseline{}{我祈求生命，也甘愿长久忍受任何劳苦。}
\verseline{}{但福尔图娜啊，倘若你正以某种不可言说的灾祸相威胁，}
\verseline{}{那么现在，啊，就在现在，让我结束这残酷的生命；}
\verseline{580}{趁忧虑仍悬而未决，未来的希望仍未分明，}
\verseline{}{趁我仍在怀中拥着你，亲爱的孩子，}
\verseline{}{我迟暮而唯一的欢喜——别让更沉重的消息}
\verseline{}{刺伤我的耳朵。”父亲在最后一次离别中倾吐这些话；}
\verseline{}{随即倒下，由仆人抬回屋内。}
\verseline{585}{这时骑兵已经从敞开的城门出发，}
\verseline{}{埃涅阿斯和忠实的阿卡忒斯行在最前，}
\verseline{}{其后是其他特洛伊贵族；帕拉斯本人处在队伍}
\verseline{}{中央，斗篷和绘彩的武装使他分外醒目：}
\verseline{}{如同晨星从大洋的水中沐浴而出，\footnote{帕拉斯在队伍中央如晨星从大洋升起；晨星是维纳斯最爱的星辰。帕拉斯披着明亮斗篷和绘彩武装出发，紧接着又被比作冲破黑暗的晨星。}}
\verseline{590}{在所有星辰之火中，维纳斯最爱这一颗；}
\verseline{}{它把神圣的面容举上天空，解开黑暗。}
\verseline{}{惶恐的母亲们立在城墙上，双眼追随着}
\verseline{}{尘土形成的云与青铜闪耀的队列。}
\verseline{}{骑兵全副武装，穿过灌木，沿最直接的路程}
\verseline{595}{向前推进；呼声传开，队形已经列成，}
\verseline{}{四蹄齐响，以重声敲击松软的田野。}
\verseline{}{凯雷清冷的河流附近，有一片巨大的圣林，\footnote{凯雷河附近的圣林奉献给田野与牲畜之神西尔瓦努斯，相传最早由佩拉斯吉人设立祭日。塔耳孔与伊特鲁里亚军队驻扎不远；埃涅阿斯先在这里与盟军会合。}}
\verseline{}{因祖先留下的敬畏而远近皆被奉为圣地；四周}
\verseline{}{山丘将它环抱，黑色冷杉围住林苑。}
\verseline{600}{相传最早占据拉丁疆界的古老佩拉斯吉人}
\verseline{}{把这片林苑和祭日奉献给西尔瓦努斯，}
\verseline{}{那位田野与牲畜的神。}
\verseline{}{离这里不远，塔耳孔与伊特鲁里亚人占据}
\verseline{}{设防稳固的营地；从高耸的山丘望去，}
\verseline{605}{整支军团已可尽收眼底，在辽阔田野上铺展。}
\verseline{}{父亲埃涅阿斯与为战争选出的青年走入营地，}
\verseline{}{安顿疲惫的马匹，也让自己的身体休息。}
\verseline{}{这时，洁白的女神维纳斯穿过天上的云，\footnote{库忒瑞亚是维纳斯得自库忒拉岛的称号。她交付的胸甲呈血红色，如青蓝云层被阳光点燃；胫甲中的金银合金与精炼黄金也属于材料描写。诗刚称盾面构图“不可尽述”，随即用百余行逐项叙写。}}
\verseline{}{带着赠礼来到。她远远看见儿子}
\verseline{610}{独自在隐僻的山谷中，靠近一条清冷的河，}
\verseline{}{便主动显现，对他说出这些话：}
\verseline{}{“看啊，我丈夫以自己的技艺完成了许诺的赠礼；}
\verseline{}{儿子，此后你不必迟疑，尽可向傲慢的}
\verseline{}{劳伦图姆人或凶猛的图尔努斯求战。”}
\verseline{615}{库忒瑞亚说完，请求儿子的拥抱，}
\verseline{}{又把放光的武具置于对面一棵橡树下。}
\verseline{}{他为女神的赠礼和如此重大的荣誉欢喜，}
\verseline{}{目光怎么也看不满足，逐件流转；}
\verseline{}{他惊叹，把一件件武具放在手掌与手臂间转动：}
\verseline{620}{那顶饰有盔鬃、吐出火焰的可怖头盔，}
\verseline{}{带来死亡的剑，僵硬的青铜胸甲，}
\verseline{}{巨大、血红——如同青蓝的云}
\verseline{}{被太阳光线点燃，远远放出辉芒；}
\verseline{}{还有金银合金与精炼黄金制成的光滑胫甲，}
\verseline{625}{长枪，以及那不可尽述的盾面构图。}
\verseline{}{火之主把意大利的事迹、罗马人的凯旋\footnote{维纳斯请武尔坎努斯为埃涅阿斯打造甲胄，沿用《伊利亚特》中忒提斯请赫淮斯托斯为阿喀琉斯造甲的场面。荷马的盾面铺开人间与宇宙，维吉尔则让盾面出现尚未发生的罗马历史；埃涅阿斯看见这些图像，却不知道它们所指的未来。}}
\verseline{}{铸在盾上；他深知预言，也深知未来的时代；}
\verseline{}{盾上还有从阿斯卡尼乌斯延续而来的未来}
\verseline{}{全部族系，以及依次发生的战争。}
\verseline{630}{他铸出一头刚刚产子的母狼，卧在青绿的}
\verseline{}{玛尔斯洞穴中；双生的男孩悬在她的乳房旁}
\verseline{}{嬉戏，无所畏惧地舔着母亲；}
\verseline{}{母狼回转圆润的颈项，轮流抚爱两个孩子，}
\verseline{}{以舌头舔塑他们的身体。\footnote{母狼与双生子对应罗慕路斯、雷穆斯的起源传说。盾面上的母狼刚刚产子，回转颈项轮流抚爱两个孩子，以舌头舔塑他们的身体；双生子则在乳房旁嬉戏、舔舐母亲。}}
\verseline{635}{离这里不远，是罗马和不顾礼法被抢走的萨宾妇女：\footnote{罗慕路斯一方在竞技集会中抢走萨宾妇女，引发与提图斯·塔提乌斯及库瑞斯人的战争。盾面随后立刻转向和约：两位国王持祭酒盘站在朱庇特坛前，杀母猪结盟；相邻图像先写抢夺与战争，随即写祭牲和结盟。}}
\verseline{}{宏大的竞技举行之际，她们正坐在看台席上；}
\verseline{}{随后一场崭新的战争突然兴起，}
\verseline{}{一边是罗慕路斯的子民，一边是年老的塔提乌斯和严厉的库瑞斯人。}
\verseline{}{战事放下以后，两位国王全副武装，}
\verseline{640}{站在朱庇特祭坛前，手持祭酒盘，}
\verseline{}{杀死一头母猪，以血缔结合约。}
\verseline{}{离那里不远，疾驰的四马战车向两边奔离，}
\verseline{}{把墨提乌斯撕开——而你，阿尔巴人啊，本该信守诺言；\footnote{阿尔巴统帅墨提乌斯曾允诺援助罗马，却在战斗中观望胜负；图卢斯·霍斯提利乌斯以背信为由，把他缚在两辆向相反方向疾驰的战车上撕裂。诗人直接呼告墨提乌斯“本该信守诺言”，又把处刑写得极其具体：内脏被拖过树林，鲜血滴在灌木上。}}
\verseline{}{图卢斯拖着那失信者的内脏穿过树林，}
\verseline{645}{血液四处洒落，从灌木上滴下。}
\verseline{}{波耳塞纳又命罗马人接纳被逐的塔耳奎尼乌斯，}
\verseline{}{以一场巨大的围城压迫罗马；}
\verseline{}{埃涅阿斯的后裔为自由奔向刀剑。}
\verseline{}{你会看见波耳塞纳满怀愤怒，仿佛正在威胁，}
\verseline{650}{因为科克勒斯竟敢拆毁桥梁，}
\verseline{}{克洛厄莉娅也挣断束缚，游过河流。\footnote{伊特鲁里亚王波耳塞纳围困罗马，要求恢复被逐的塔耳奎尼乌斯王政。科克勒斯守桥并令同伴毁桥，克洛厄莉娅则挣脱人质束缚、游过台伯河；两人的行动在盾面上共同代表罗马人为自由抵抗王权复辟。}}
\verseline{}{在盾面最高处，曼利乌斯守卫塔耳佩亚城堡，}
\verseline{}{立在神庙前，占据高耸的卡庇托利乌姆；}
\verseline{}{罗慕路斯新建的宫殿仍覆着粗硬的茅草。}
\verseline{655}{这里，一只银鹅在金色的柱廊间振翅飞过，}
\verseline{}{以叫声警告高卢人已经来到门槛前。}
\verseline{}{高卢人穿过灌木，已经逼近城堡，}
\verseline{}{受黑暗和昏沉夜色的赠礼掩护；}
\verseline{}{他们头发金黄，衣装也呈金色，}
\verseline{660}{条纹短斗篷闪亮，乳白的颈项}
\verseline{}{缠着黄金；每个人手中挥舞两支}
\verseline{}{阿尔卑斯长矛，身体由狭长的盾牌遮护。}
\verseline{}{这里，武尔坎努斯又锤出欢跃的萨利祭司、赤裸的卢佩耳库斯祭司，\footnote{曼利乌斯守卫卡庇托利乌姆，朱诺的圣鹅以叫声警告高卢人夜袭。随后出现萨利祭司、近乎赤裸的卢佩耳库斯祭司、羊毛尖帽、从天而降的圣盾和贵妇祭车；盾面在早期战争图像之后又刻出后世罗马的祭司、圣物与祭车。}}
\verseline{}{头戴羊毛尖帽的人，以及从天而降的圣盾；}
\verseline{665}{贞洁的贵妇们坐在柔软的祭车中，}
\verseline{}{引领圣物穿过城中。远处，他又加上}
\verseline{}{塔耳塔洛斯的居地和狄斯深邃的门，}
\verseline{}{罪行所受的惩罚；还有你，喀提林，}
\verseline{}{悬在威胁的峭壁上，面对复仇女神的面孔战栗；}
\verseline{670}{远离他的另一处是尽责者，加图在他们中间颁布法则。\footnote{喀提林被置于塔耳塔洛斯受刑，复仇女神在近旁威胁；尽责者与他分隔，加图为他们颁布法则。塞尔维乌斯（8.670）解释这里的加图为老加图；另有小加图解（John Conington，\href{https://www.perseus.tufts.edu/hopper/text?doc=Perseus:text:1999.02.0051:book=8\&force=y}{《埃涅阿斯纪评注》，评8.670，Perseus 数字本}），或认为诗人有意兼指两者。}}
\verseline{}{在这些景象之间，涨满的海洋图像辽阔铺开，}
\verseline{}{通体金黄，青蓝海水却翻卷着白色泡沫；}
\verseline{}{闪亮的银海豚环绕一周，以尾巴}
\verseline{}{扫过海面，切开潮流。}
\verseline{675}{正中央可以看见青铜闪耀的舰队和阿克提乌姆战事；\footnote{盾面先铺开金色海洋，写到青蓝海水、白色泡沫和银海豚；阿克提乌姆舰队随后出现在海面中央。}}
\verseline{}{战阵一旦排开，整个琉卡忒斯仿佛沸腾，}
\verseline{}{海浪也因黄金灿然发光。}
\verseline{}{这边，奥古斯都·恺撒率意大利人走向战斗，\footnote{盾面把意大利人、元老、人民、家神与大神全部列在奥古斯都·恺撒一方。头顶的“父祖星辰”通常指神化尤利乌斯·恺撒的彗星；阿格里帕则佩饰有舰首的海战冠，这是表彰海战功绩的荣标。}}
\verseline{}{与他同在的有元老和人民、家神与伟大的众神；}
\verseline{680}{他立在高高的船尾，欢然的双鬓}
\verseline{}{吐出两道火焰，头顶显出父祖的星辰。}
\verseline{}{另一处，阿格里帕得顺风与众神相助，}
\verseline{}{身形高峻，率领军阵；他的额头上——骄傲的}
\verseline{}{战争荣标——闪着饰有舰首的海战冠。}
\verseline{685}{那边，安东尼倚仗异邦的力量与各式武装，\footnote{安东尼被配置在东方、埃及、巴克特里亚和“异邦力量”一侧；诗中不直呼克娄巴特拉之名，只称她为 Aegyptia coniunx（“埃及妻子”）。插入句中的 nefas 是叙述者的惊呼，表示违逆神律、可憎可耻之事；这里尤其谴责一位罗马统帅带着埃及女王和东方军力投入罗马内战。}}
\verseline{}{从奥罗拉的民族和红海沿岸凯旋而来，}
\verseline{}{带着埃及、东方的军力和遥远的巴克特里亚；}
\verseline{}{他的埃及妻子竟也跟随其后——大逆不道！}
\verseline{}{所有舰船一齐冲向前，整个海面被后掠的}
\verseline{690}{桨叶和三叉舰首撕碎，翻卷白沫。}
\verseline{}{他们驶向深海；你会以为拔根而起的库克拉得斯群岛}
\verseline{}{正在海上漂浮，又以为一座座高山相互撞击——}
\verseline{}{男子们站在筑有塔楼的舰船上，以如此巨大的体量进逼。}
\verseline{}{手中掷出燃着麻絮的飞铁投矢，}
\verseline{695}{四处洒落；涅普顿的田野被新近的屠杀染红。}
\verseline{}{女王立在军阵中央，以祖国的叉铃召集队伍，}
\verseline{}{还没有看见背后那一双毒蛇。\footnote{女王以埃及叉铃召集军阵，仍未看见身后双蛇；双蛇在盾面上预示她战败后的自杀。盾面把犬首阿努比斯与“各式神怪”放在罗马诸神的敌阵一侧。女王逃离后，尼罗河却展开衣襟，哀恸地召唤败者进入青蓝怀抱与隐秘水道。}}
\verseline{}{各种各样的神怪和犬首吠叫的阿努比斯，}
\verseline{}{举起兵器，对抗涅普顿、维纳斯与弥涅耳瓦。}
\verseline{700}{玛沃耳斯在战斗中央肆虐，形象被刻入钢铁；}
\verseline{}{阴沉的复仇女神也从高空降临，}
\verseline{}{不和女神披着撕裂的斗篷，欢喜地向前行进，}
\verseline{}{贝罗娜挥着染血的长鞭紧随其后。}
\verseline{}{阿克提乌姆的阿波罗从高处看见这一切，正在拉弓；}
\verseline{705}{被这恐怖驱赶，整个埃及、印度人、}
\verseline{}{整个阿拉伯和所有萨巴人都转身逃跑。}
\verseline{}{女王本人仿佛已经召来风息，张开}
\verseline{}{船帆，眼看就要放松帆索；}
\verseline{}{火之主铸出她穿过屠杀，因将临的死亡而苍白，}
\verseline{710}{被海浪和伊阿庇癸风带走。}
\verseline{}{另一边，尼罗河以巨大的身躯哀恸，}
\verseline{}{展开衣襟，以整件长袍呼唤战败者，}
\verseline{}{进入青蓝的怀抱和隐秘的水道。}
\verseline{}{而恺撒乘着三重凯旋进入罗马城内，\footnote{“三重凯旋”联系奥古斯都在公元前 29 年为伊利里亚、阿克提乌姆与埃及举行的连续凯旋。三百座神庙、全城欢声、妇女合唱、祭坛与被宰公牛都出现在凯旋场面中；奥古斯都随后坐在福玻斯神庙的雪白门槛，接受战败诸族的礼物。}}
\verseline{715}{向意大利诸神献上不朽的誓愿——}
\verseline{}{三百座宏伟神庙，遍布全城。}
\verseline{}{道路因欢乐、竞技与掌声轰鸣；}
\verseline{}{所有神庙都有妇女合唱，所有祭坛都在举行祭礼；}
\verseline{}{被宰杀的公牛铺满祭坛前的地面。}
\verseline{720}{奥古斯都本人坐在灿烂福玻斯雪白的门槛上，}
\verseline{}{检阅各民族的贡礼，把它们悬上傲然的}
\verseline{}{门柱；战败的民族排成长列走来，}
\verseline{}{语言彼此不同，衣袍的穿法和兵器也各不相同。}
\verseline{}{这里，穆尔喀伯铸出诺玛德人和衣带松散的阿非利人，}
\verseline{725}{铸出勒勒格斯人、卡里亚人与持箭的革洛诺伊人；}
\verseline{}{幼发拉底河走来，水势已变得更加温顺，}
\verseline{}{还有居于人类尽头的摩里尼人、双角的莱茵河、}
\verseline{}{不肯屈服的达哈伊人，以及愤恨桥梁的阿拉克斯河。\footnote{凯旋长列包括北非的诺玛德人和阿非利人，以及勒勒格斯人、卡里亚人、革洛诺伊人、摩里尼人和达哈伊人。幼发拉底河、双角莱茵河与阿拉克斯河也被人格化，同诸族一起列入凯旋图像，勾勒诗中罗马征服所及的地域。幼发拉底河水势变得温顺，阿拉克斯河却仍愤恨桥梁。}}
\verseline{}{埃涅阿斯惊叹武尔坎努斯盾上的这一切——母亲的赠礼；}
\verseline{730}{他并不知晓那些事，却因图像而欣喜，}
\verseline{}{把后裔的声名与命运举上肩头。\footnote{“母亲的赠礼”指维纳斯带来的武尔坎努斯之盾。埃涅阿斯因图像欣喜，却“并不知晓”图像所指事件；最后他把盾举上肩，也把写在盾上的后裔声名与命运一并背起。}}
